\chapter{相与相变}
\label{chap:phase}

\begin{definition}
	{相}{phase}
	一个相(phase)是系统中组成和物理性质在宏观上均匀的部分。例如晶体的微观密度可随晶格周期变化，宏观取平均后仍属于同一相。
	比如：
	\begin{itemize}
		\item 晶态固体(crystalline solid)：原子、离子或分子在晶格上排列，具有长程位置有序性；
		\item 液相(liquid phase)：分子在空间上无序排列，具有短程有序性；
		\item 气相(gas phase)：分子没有长程位置有序性，在稀薄极限下分子间作用对状态方程的影响很小。
	\end{itemize}
\end{definition}

\begin{remark}
	一个相可以是单组分，例如纯液态水，也可以包含多个组分，例如均匀溶液。相内的物种还可以发生化学反应并达到反应平衡。
\end{remark}

% 在一定条件下，不同的相可以共存。不同相之间的转变称为相变。相变发生在外部条件发生变化时，即强度量变化时。相变是竞争的结果。微观上有各种不同的相互作用。不同的相互作用可能会导致不同类型的相变。系统中两个能量尺度间的关系在强度量变化时发生变化。这种关系变化到一定程度后，系统出现新的相。例如，微观上原子分子的相互作用和热运动的相互竞争。控制这两个因素的参量可以是温度（也可以是压强）。在低温时，物质倾向于形成有长程关联的有序结构，这是分子间相互作用的结果。在高温，熵的作用占优势。这种竞争导致相变。相变可以发生在某一温度处，即相变温度，该温度将低温相和高温相隔开。在降温过程中，当微观上相互作用的能量与热运动能量可比拟时，相变发生，出现新的秩序。

\section{化学势}
\label{sec:chemical potential}

我们在前文没有讨论开放系统中的平衡问题。
与封闭系统不同，开放系统的摩尔数$n$不是恒定不变的，由此需要引入新的强度量。

\begin{definition}
	{化学势}{chemical potential}
	在确定的温度$T$和压强$p$下，系统的化学势是Gibbs自由能对摩尔数的偏导数：
	\begin{equation}
		\mu=\pvat{G}{n}{T,p}.
	\end{equation}
\end{definition}

\begin{corollary}
	利用Euler公式可得$G=n\mu$，摩尔Gibbs自由能$g=\mu$就是化学势。
\end{corollary}

\begin{corollary}
	内能的微分式、Euler方程和Gibbs--Duhem关系分别为
	\begin{subequations}
		\begin{gather}
			\d U=T\d S-p\d V+\mu\d n,\\
			U=TS-pV+\mu n,\\
			S\d T-V\d p+n\d\mu=0.
		\end{gather}
	\end{subequations}
	化学势的引入也带来了更多的Maxwell关系。
\end{corollary}

\begin{example}
	{理想气体的化学势}{}
	由$\mu=G/n$可得
	\begin{equation}
		\mu=\mu_0\frac T{T_0}-RT\ln\biggb{\biggp{\frac T{T_0}}^{c+1}\frac{p_0}p}.
	\end{equation}
\end{example}

\section{相变}
\label{sec:phase transition}

考虑单组份的孤立系统中共存的两个相(记为$\phaseI,\phaseII$)，有
\begin{align*}
	U_\phaseI+U_\phaseII&=U_\tot,\\
	V_\phaseI+V_\phaseII&=V_\tot,\\
	n_\phaseI+n_\phaseII&=n_\tot.
\end{align*}
平衡条件要求
\[
	\d S=\d S_\phaseI+\d S_\phaseII
	=\biggp{\frac1{T_\phaseI}-\frac1{T_\phaseII}}\d U_\phaseI
	+\biggp{\frac{p_\phaseI}{T_\phaseI}-\frac{p_\phaseII}{T_\phaseII}}\d V_\phaseI
	-\biggp{\frac{\mu_\phaseI}{T_\phaseI}-\frac{\mu_\phaseII}{T_\phaseII}}\d n_\phaseI
	=0.
\]
可得
\begin{subequations}
	\begin{align}
		T_\phaseI&=T_\phaseII,\\
		p_\phaseI&=p_\phaseII,\\
		\mu_\phaseI&=\mu_\phaseII.
	\end{align}
\end{subequations}
因此，忽略界面和外力场的影响时，两相平衡要求温度、压强和化学势分别相等。在温度和压强相同的条件下，物质从化学势高的一相转入化学势低的一相会降低总Gibbs自由能，转变可持续到两相达到平衡或其中一相耗尽。

考虑相变的方向，物体最常见的三种相——固、液、气——之间的转变共有6种。每种都有对应的名称，如\figref{fig:phase transition} 所示。
这些相变的过程在日常生活中也是常见的。

\begin{notecaptionblock}
	\tikzfigure{phase-transitions}
	\captionof{figure}{固液气相变}
	\label{fig:phase transition}
\end{notecaptionblock}

\subsection{相图}
\label{ssec:phase graph}

% 显然，不同温度$T$、压强$p$和体积$V$下，
表示相平衡系统的组成与一些参数（如$p,T,V,\mu$）之间关系的图叫相图。
% 同一温度$T$、压强不同相都有各自的化学势$\mu_{\mathrm{X}}$，但只有化学势最低的相可以稳定存在。
% 可想而知，已知的各个相都对应着相图中的一块区域。

\begin{notecaptionblock}
	\tikzfigure{water-pvt}
	\captionof{figure}{水在温度 $T$、体积 $V$ 和压强 $p$ 构成的三维空间中的相图}
	\label{fig:H2O phase P-T-V}
	% \footnote{这张图花了我老长时间了。}
\end{notecaptionblock}

\paragraph{$p\vs T$相图}

将\figref{fig:H2O phase P-T-V} 投影到$p\vs T$坐标上，如\figref{fig:H2O phase P-T} 所示，
% 对于单一组分而言，某一相的化学势$\mu(T,p)$只是温度$T$和压强$p$的函数，
固液气三相分别占据了不同的区域，区域的边界线上两个相共存。
沿共存线，Clapeyron 方程给出熔点、沸点或升华点随压强的变化率 $\d T/\nd p$。在线性压强轴上曲线斜率为 $\d p/\nd T$，本图采用对数压强轴。

\begin{theorem}
	{Clapeyron方程}{Clapeyron equation}
	沿两相共存曲线有$\mu_\phaseI=\mu_\phaseII$，
	由$\d\mu=-s\d T+v\d p$可得
	\[
		(s_\phaseI-s_\phaseII)\d T=(v_\phaseI-v_\phaseII)\d p\implies\dv Tp=\frac{v_\phaseI-v_\phaseII}{s_\phaseI-s_\phaseII}.
	\]
	一摩尔物质由$\phaseI$相转为$\phaseII$相时吸收的相变潜热(latent heat)为$\ell:=T(s_\phaseII-s_\phaseI)$，故
	\begin{equation}
		\dv Tp=\frac{T(v_\phaseII-v_\phaseI)}\ell.
	\end{equation}
\end{theorem}

\begin{corollary}
	取$\ell$为摩尔汽化潜热。远离临界点时，若$v_\gas\gg v_\liquid$，气相可视为理想气体，且在所取温区内可将$\ell$视为常数，则
	\[
		\dv Tp\approx\frac{Tv_\gas}\ell=\frac{RT^2}{\ell p}\implies p=p_0\exp\biggb{\frac\ell{R}\biggp{\frac1{T_0}-\frac1T}}.
	\]
	
\end{corollary}

\begin{remark}
	常压附近水结成普通冰时存在反常膨胀，$v_\liquid<v_\solid$，因此这段熔化曲线的$\d T/\d p<0$。
	% 在Celsius温标下，$\qty{1}{\atm}=\qty{101325}{Pa}$下水的凝固点和沸点分别为$\qty{0}{\degreeCelsius}$和$\qty{100}{\degreeCelsius}$。
\end{remark}

\begin{notecaptionblock}
	\tikzfigure{water-pt}
	\captionof{figure}{水的$p\vs T$相图(示意图)}
	\label{fig:H2O phase P-T}
\end{notecaptionblock}

特别地，随着温度的升高，气液共存线会存在一个终点，称为临界点(critical point)。
当$T>T_\criticalpoint$时，流体可以从低密度连续压缩到高密度，途中不再跨越气液相界。
这将在\secref{ssec:liquid-gas phase transition}讨论。
水的临界点$T_\criticalpoint=\qty{647.1}{K},p_\criticalpoint=\qty{22.064}{MPa}$。

三条共存线的交点称为三相点(triple point)，此处固、液、气三相共存。
水的三相点约为$T_\triplepoint=\qty{273.16}{K},p_\triplepoint=\qty{611.7}{Pa}$。

\paragraph{$p\vs V$相图}

投影到 $p\vs V$ 平面后可画出等温线，如\figref{fig:H2O phase P-V} 所示。
在图示的流体区域内，$T>T_\criticalpoint$的等温线全部位于单相区；
三相点与临界点之间的等温线穿过气液共存区。给定温度，气液共存压强便已确定，因此共存段是一条水平线。

\begin{example}
	{杠杆原理}{}
	气液共存区中的等温线代表确定的$p,T$下可以有不同的体积，这与气液两相的摩尔比有关。
	等温线在气液共存区的两端表示全为液相或气相，对应的摩尔体积为$v_\liquid,v_\gas$；
	而在共存区内，若液、气相所占的摩尔比分别为$x_\liquid,x_\gas$，则其摩尔体积满足杠杆原理：
	\[
		v=v_\liquid x_\liquid+v_\gas x_\gas\iff x_\gas=\frac{v-v_\liquid}{v_\gas-v_\liquid}.
	\]
\end{example}

\begin{notecaptionblock}
	\tikzfigure{pv-isotherms}
	\captionof{figure}{$p\vs V$相图(示意图)和等温线}
	\label{fig:H2O phase P-V}
\end{notecaptionblock}

对于通常的液气临界点，等温压缩率发散，等温线在临界点C的斜率$(\p p/\p V)_T$趋于零。稳定单相支满足
\[
	\pvat{p}{V}{T}=-\frac{1}{V\kappa_T}<0 .
\]
若临界等温线在C附近二阶可微，一阶导数在C处达到局部最大值零，从而有水平拐点条件
\begin{equation}
	\label{eqn:critical point}
	\pvat{p}{V}{T}=0,\quad\biggp{\pv[2]pV}_T=0.
\end{equation}

\begin{example}
	{van der Waals气体的约化变量}{vdW critical}
	理想气体不存在液化，考虑van der Waals气体。
	% 对于临界温度$T_\criticalpoint$
	% \begin{align*}
	% 	\pv pv&=-\frac{RT}{(v-b)^2}+\frac{2a}{v^3}=0, \\
	% 	\pv[2]pv&=\frac{2RT}{(v-b)^3}-\frac{6a}{v^4}=0.
	% \end{align*}
	代入临界点的拐点条件，解得
	\begin{subequations}
		\label{eqn:vdW critical}
		\begin{align}
			v_\criticalpoint&=3b,\\
			T_\criticalpoint&=\frac{8}{27}\frac{a}{Rb},\\
			p_\criticalpoint&=\frac{RT_\criticalpoint}{v_\criticalpoint-b}-\frac a{v_\criticalpoint^2}=\frac1{27}\frac a{b^2}.
		\end{align}
	\end{subequations}
	可以定义约化变量$\tilde T:=T/T_\criticalpoint$等，可得状态方程：
	\begin{equation}
		\tilde p=\frac{8\tilde T}{3\tilde v-1}-\frac3{\tilde v^2}.
	\end{equation}
\end{example}

\paragraph{$T\vs V$相图}

投影到$T\vs V$图上与$p\vs V$图类似。
% 只有$p<p_\criticalpoint$时才会有气液共存区，共存区的等压线垂直于$T$轴（对应沸点）。

\subsection{气液相变}
\label{ssec:liquid-gas phase transition}

固定摩尔数$n$时，每个相的Gibbs自由能都是$p,T$的连续函数。平衡态取各相中最低的$G$，两相共存时二者相等。跨越一级相变边界时，平衡的$G$连续而斜率突变，曲面像一张有折痕的纸，如\figref{fig:phase G-p-T} 所示。

\begin{notecaptionblock}
	\tikzfigure{gibbs-pt}
	\captionof{figure}{气液两相的 Gibbs 自由能曲面 $G(p,T)$}
	\label{fig:phase G-p-T}
\end{notecaptionblock}

\begin{definition}
	{一级相变和连续相变}{}
	至少有一个Gibbs自由能的一阶导数不连续的相变称为一级相变，否则称为连续相变。
\end{definition}

\begin{corollary}
	连续相变的摩尔熵$s=-(\p g/\p T)_p$是连续的，因此没有潜热。
\end{corollary}

\iffalse
\begin{notecaptionblock}
	\tikzfigure{transition-gibbs-energy}
	\tikzfigure{transition-entropy}
	\tikzfigure{transition-volume}
	\captionof{figure}{相变点处Gibbs自由能}
\end{notecaptionblock}
\fi

\paragraph{共存压强与Maxwell等面积法则}

在低于临界温度的等温线上压缩气体，进入共存区后，压缩使部分蒸气液化，压强保持不变。为求出平台的压强，设活塞下有$n$摩尔流体，温度$T$和外压$p_{\rm ext}$给定。以$v=V/n$为试取的摩尔体积，$f(T,v)$为均匀流体的摩尔Helmholtz自由能。系统膨胀还须对外做功，故应比较
\[
	\Phi(v)=f(T,v)+p_{\rm ext}v,\qquad
	\Phi'(v)=p_{\rm ext}-p(T,v),\qquad
	\Phi''(v)=-\pvat{p}{v}{T}.
\]
驻点满足内压等于外压，此时$\Phi=g=\mu$。在\secref{sec:thermo-stability}的力学稳定支上，$(\partial p/\partial v)_T<0$，故$\Phi$取局部极小值。平衡时选取全局极小值。

以van der Waals流体为例，由$(\partial f/\partial v)_T=-p$沿等温线积分得
\[
	p(T,v)=\frac{RT}{v-b}-\frac a{v^2},\qquad
	f(T,v)=f_0(T)-RT\ln\frac{v-b}{b}-\frac av,\quad v>b .
\]
当$T<T_\criticalpoint$时，等温线可与同一水平线交于三点。若外侧两个交点$v_\liquid,v_\gas$处的$\Phi$相等，则液相与气相具有同一化学势。记此时的外压为$p_{\rm sat}(T)$，由$\Phi(v_\liquid)=\Phi(v_\gas)$可得
\[
	\frac{f(T,v_\gas)-f(T,v_\liquid)}{v_\gas-v_\liquid}
	=-p_{\rm sat}
	=\left.\pv f v\right|_{T,v_\liquid}
	=\left.\pv f v\right|_{T,v_\gas}.
\]
连接两点的直线是$f(T,v)$的公切线，沿这条直线的摩尔自由能由两相按杠杆原理混合得到。用$f(T,v_\gas)-f(T,v_\liquid)=-\int_{v_\liquid}^{v_\gas}p(T,v)\d v$消去自由能，即得\textbf{Maxwell等面积法则}
\begin{equation}
	\int_{v_\liquid}^{v_\gas}\bigb{p(T,v)-p_{\rm sat}}\d v=0,\qquad
	p(T,v_\liquid)=p(T,v_\gas)=p_{\rm sat}.
\end{equation}
例如取$\tilde T=0.90$。对于试取的$\tilde p_{\rm sat}$，从约化等温线的三个交点中取最小、最大的根$\tilde v_\liquid,\tilde v_\gas$。等面积条件写为
\[
	\tilde p_{\rm sat}(\tilde v_\gas-\tilde v_\liquid)
	=\frac{8\tilde T}{3}\ln\frac{3\tilde v_\gas-1}{3\tilde v_\liquid-1}
	+3\biggp{\frac1{\tilde v_\gas}-\frac1{\tilde v_\liquid}} .
\]
调节共存压强使上式成立，得到$\tilde p_{\rm sat}\simeq0.6470$、$\tilde v_\liquid\simeq0.6034$和$\tilde v_\gas\simeq2.3488$。在这两个体积之间，均匀支的两个转折点分别位于$\tilde v\simeq0.7186,1.5285$。两转折点之间的均匀流体有$\kappa_T<0$，对微小体积涨落已经失稳。从各共存端点延伸到相邻转折点的均匀支仍有$\kappa_T>0$，但两相混合的自由能更低，这一段是亚稳支。

\begin{notecaptionblock}
	\tikzfigure{gibbs-phase-fraction}
	\captionof{figure}{$T=0.90T_\criticalpoint$的van der Waals等温线。水平线两端的实心点给出共存体积，两个阴影区域的面积相等。空心点标出均匀支的力学失稳边界。}
	\label{fig:model of liquid-gas phase transition}
\end{notecaptionblock}

\paragraph{过热与成核}

在给定外压下将液体加热到沸点以上时，气相的化学势较低。只要液相的局部极小值还存在，液体就可能暂时保持均匀，这种状态称为过热(overheating)。\figref{fig:overheating}中的两个极小值分别代表同一温度和外压下的均匀液相与气相。
\begin{notecaptionblock}
	\tikzfigure{superheated-liquid}
	\captionof{figure}{$T=0.90T_\criticalpoint$、$p_{\rm ext}=0.63p_\criticalpoint$时均匀流体的摩尔势$\Phi$，以气相极小值$\Phi_\gas$为参考。液相位于较高的局部极小值。}
	\label{fig:overheating}
\end{notecaptionblock}
实际相变往往从局部热涨落产生的新相小核开始，称为成核(nucleation)。以蒸气凝结为例，在给定压强下将蒸气冷却至共存温度$T_V$以下，设过冷度$\Delta T=T_V-T>0$。形成液滴会增加界面能，液滴内部的物质转为化学势更低的液相则会降低自由能。设液滴为半径$R$的球，液相摩尔密度$\tilde n_\liquid$和表面张力$\sigma$近似不随半径改变。两项贡献给出形成液滴所需的自由能
\begin{equation}
	\Delta G(R)=4\pi R^2\sigma-\frac{4\pi}3 R^3\tilde n_\liquid\Delta\mu,\qquad
	\Delta\mu=\mu_\gas(T,p)-\mu_\liquid(T,p)>0.
\end{equation}
在共存点附近，$(\partial\mu/\partial T)_p=-s$给出$\Delta\mu\simeq(s_\gas-s_\liquid)\Delta T=\ell\Delta T/T_V$，其中$\ell$为摩尔汽化潜热。$\Delta G(R)$的最大值给出临界半径及成核势垒
\[
	R_\criticalpoint=\frac{2\sigma}{\tilde n_\liquid\Delta\mu}
	\simeq\frac{2\sigma T_V}{\tilde n_\liquid\ell\Delta T},\qquad
	\Delta G(R_\criticalpoint)=\frac{4\pi}3\sigma R_\criticalpoint^2 .
\]
小于临界半径的液滴倾向于缩小，较大的液滴则倾向于长大。杂质和容器壁可降低形成新界面的代价，使成核更容易发生。

\begin{notecaptionblock}
	\tikzfigure{nucleation-barrier}
	\captionof{figure}{成核临界半径$R_\criticalpoint$}
	\label{fig:nucleation}
\end{notecaptionblock}

成核现象普遍存在于一级相变的系统中，如Wilson云室、人工降雨等。


% \subsectionstar{铁磁相变}

% 在无外场的条件下，固体中具有磁矩的分子会在热运动中无序排列，对外不显磁性。
% 但铁磁物质在低于临界温度(称为Curie温度) $T_\criticalpoint$时会自发地产生磁性，形成磁有序结构，称为铁磁相变。\footnote{历史上是反过来的：磁铁在超过Curie温度时会失去磁性。}

% 我们只考虑磁矩仅取上下两个方向的情形。对于一个无外磁场的系统，上下应该等价，系统的运动方程关于上下两方向间的变换应该保持不变，即具有对上下变换的对称性。但是，在居里温度下，铁磁系统只会选择上下其中的一个方向。这个现象称为自发对称破缺。发生破缺的是系统的状态而不是系统的运动方程。温度是系统的一个强度量。铁磁系统会与外磁场相互作用。因此另一个强度量为外磁场。表示系统性质的量为磁化强度。
% 系统的相图如图10.18所示。M作为T和B的函数，如同撕开一个缝的一张纸。缝隙垂直于B-T平面，并关于M对称。

\subsectionstar{Landau连续相变理论}

\begin{definition}
	{序参量}{order parameter}
	在超导、磁性等一大类相变中，有一个区别不同相的热力学量，称为序参量(order parameter) $\eta$。
	无序相$\eta=0$；有序相$\eta\neq 0$，对应对称破缺。
\end{definition}

\begin{remark}
	$\eta$可以是复数（超导、超流），可以是二维和三维的，可以与空间的维数不同。
\end{remark}

\begin{example}
	{序参量}{}
	\begin{notecaptionblock}
		\captionof{table}{不同系统的相变序参量}
		\begin{tabular}{cccc}
			\toprule
			系统 & 相变 & 序参量 & 物理意义 \\
			\hline
			液--气 & 气化/液化 & $\tilde n_\liquid-\tilde n_\gas$ & 密度差 \\
			% 固 - 液 & 熔化/凝固 & $\tilde n_\solid-\tilde n_\liquid$ & 密度差 \\
			\multirow{2}{*}{磁性} & 铁磁/顺磁 & $M$ & 磁化强度 \\
			& 反铁磁/顺磁 & $M_{\mathrm{s}}$ & 子晶格磁化强度 \\
			\multirow{2}{*}{量子液体} & 超导 & \multirow{2}{*}{$\psi$} & 超导电子对波函数 \\
			 & 玻色原子超流 & & 凝聚体的宏观波函数 \\
			\bottomrule
		\end{tabular}
		\label{tab:order parameter}
	\end{notecaptionblock}
\end{example}

\begin{definition}
	{临界点与临界现象}{critical phenomenon}
	临界点(critical point)是连续相变的相变点，临界温度$T_\criticalpoint$；
	临界现象(critical phenomenon)是物质在临界点邻域的统计热力学行为。

	以铁磁相变为例，在热力学极限下令$H\to0^+$选定一个磁化方向。比热的奇异部分$C_{H,\rm sing}$、自发磁化强度$M$和磁化率$\chi$在临界温度两侧常表现出幂律行为，临界等温线也有相应的幂律。其指数称为临界指数(critical exponent)。以下列出领先幂次：
	\begin{subequations}
		\begin{align}
			C_{H,\rm sing}&\sim\begin{cases}
				A(T/T_\criticalpoint-1)^{-\alpha}, & T>T_\criticalpoint \\
				A'(1-T/T_\criticalpoint)^{-\alpha'}, & T<T_\criticalpoint
			\end{cases}\\
			M&=\begin{cases}
				0, & T>T_\criticalpoint \\
				B(1-T/T_\criticalpoint)^{\beta},\phantom{\quad} & T<T_\criticalpoint
			\end{cases}\\
			\chi&\sim\begin{cases}
				C(T/T_\criticalpoint-1)^{-\gamma}, & T>T_\criticalpoint \\
				C'(1-T/T_\criticalpoint)^{-\gamma'}, & T<T_\criticalpoint
			\end{cases}\\
			M&\sim DH^{1/\delta},\qqqquad\qquad T=T_\criticalpoint,\quad H\to0^+
		\end{align}
	\end{subequations}
	其中$\alpha,\alpha',\beta,\gamma,\gamma',\delta$为临界指数，$A,A',B,C,C',D$为临界幅度。比热还可有解析背景，幂次为零时也可能出现对数奇异性。
\end{definition}

\begin{theorem}
	{缩放关系}{scaling relations}
	设约化温度$t=T/T_\criticalpoint-1$，自由能密度的奇异部分具有标度形式
	\[
		g_{\rm sing}(t,H)=|t|^{2-\alpha}\mathcal G_\pm(H/|t|^\Delta).
	\]
	两侧采用相同的指数和不同的标度函数$\mathcal G_\pm$。在通常的幂律标度下，六个热力学临界指数满足
	\begin{subequations}
		\begin{align}
			\alpha+2\beta+\gamma&=2,\\
			\gamma&=\beta(\delta-1),\\
			\alpha'&=\alpha,\qquad\gamma'=\gamma .
		\end{align}
	\end{subequations}
\end{theorem}

对场求一次、两次导数，分别得到磁化强度和磁化率的标度幂次，故$\beta=2-\alpha-\Delta$、$\gamma=2\Delta+\alpha-2$。要使临界等温线的$M\propto H^{1/\delta}$在$t\to0$时有有限幅度，还需$\Delta=\beta\delta$。消去$\Delta$即得上述关系。

作为例子，在临界点附近按序参量展开Landau自由能，
\begin{align}
	F(T,\eta)=F_0(T)+\frac12a(T)\eta^2+\frac14b(T)\eta^4+\cdots,
\end{align}
这里考虑具有$\eta\leftrightarrow-\eta$对称性的系统，取临界点附近$b(T)>0$，展开中不含$\eta$的奇次幂。
$\eta$应使$F$取极小值：
\[
	\pv F\eta=\eta(a+b\eta^2)=0
	% \quad\pv[2]F\eta=a+3b\eta^2>0.
	\implies\eta=0,\pm\sqrt{-\frac ab}.
\]
$a>0$时，$\eta=0$是唯一实驻点且为极小值。$a<0$时还出现一对非零实驻点$\eta_\pm$，此时$\eta=0$为极大值。当$T\to T_\criticalpoint^-$时，两个有序极小值$\eta_\pm$趋于零，因此$a(T_\criticalpoint)=0$。
由$\eta_\pm$是极小值点：
\[
	\pv[2]F\eta=a+3b\eta_\pm^2=-2a>0
	\implies a<0,\enspace (T<T_\criticalpoint).
\]
取$a$在$T_\criticalpoint$附近的线性项及$b$的常数项，其中$a_0,b_0>0$，
\[
	a(T)=a_0\biggp{\frac T{T_\criticalpoint}-1},\quad b(T)=b_0.
\]
可得临界指数$\beta=1/2$：
\begin{equation}
	\eta=
	\begin{cases}
		0,&T\geqslant T_\criticalpoint\\
		\pm\sqrt{\frac{a_0}{b_0}\biggp{1-\frac T{T_\criticalpoint}}},&T<T_\criticalpoint
	\end{cases}
\end{equation}
\begin{notecaptionblock}
	\tikzfigure{landau-free-energy}
	\captionof{figure}{Landau自由能}
	\label{fig:Landau free energy}
\end{notecaptionblock}
熵是连续的：
\begin{equation}
	S=-\pv FT=S_0-\frac{a_0}{2T_\criticalpoint}\eta^2=
	\begin{cases}
		S_0,                             & T\geqslant T_\criticalpoint \\
		S_0+\frac{a_0^2}{2b_0T_\criticalpoint}\biggp{\frac T{T_\criticalpoint}-1}, & T<T_\criticalpoint
	\end{cases}
\end{equation}
比热却不连续了：
\begin{equation}
	C=T\pv ST=
	\begin{cases}
		C_0,                 & T\geqslant T_\criticalpoint \\
		C_0+\frac{a_0^2}{2b_0T_\criticalpoint^2}T, & T<T_\criticalpoint
	\end{cases}
\end{equation}
因此是\textbf{二级相变}。有序相的比热$>$无序相，且$T=T_\criticalpoint$处比热的突变是有限的，$\alpha=\alpha'=0$。

\paragraph{外加场$B$}
以磁化强度$M$为序参量，记施加的场为$B=\mu_0H$。单位体积自由能中，外场与磁矩的耦合给出一项$-BM$，
\[
	g(M,B)=g_0+\frac12aM^2+\frac14bM^4-BM.
\]
平衡时
\[
	\pv gM=aM+bM^3-B=0.
\]
$T=T_\criticalpoint$时，$a=0,\;B=bM^3$故$\delta=3$。
体磁化率为$\chi=(\partial M/\partial H)_T=\mu_0/(a+3bM^2)$。从$B\to0^+$取极限，在临界温度以下选定一个有序支$M^2=-a/b$，得到
\[
	\chi=
	\begin{cases}
		\frac{\mu_0}{a_0}\biggp{\frac T{T_\criticalpoint}-1}^{-1},&T>T_\criticalpoint \\
		\frac{\mu_0}{2a_0}\biggp{1-\frac T{T_\criticalpoint}}^{-1}, & T<T_\criticalpoint .
	\end{cases}
\]
$\gamma'=\gamma=1$。这里用同一个$M$描述整个样品。要检验这个近似，就要知道一个区域的磁化涨落有多大，以及它和周围多大范围内的涨落相关。

\subsection{临界涨落与重整化群}
\label{sec:critical-fluctuations}

\paragraph{从自由能到空间关联}
在$d$维样品中，以$\phi(\bm x)$表示位置$\bm x$附近小区域内的平均序参量。区域的线度大于微观作用距离，又小于待研究的涨落尺度。先考虑相互作用各向同性、短程，且零场下具有$\phi\leftrightarrow-\phi$对称性的标量系统。将均匀Landau自由能推广为
\begin{equation}
    \mathcal F[\phi]=\int\d^d x\,
    \left[\frac{\kappa}{2}|\nabla\phi|^2+
    \frac{a_0t}{2}\phi^2+\frac{u}{4}\phi^4-h\phi\right],
    \qquad t=\frac{T-T_\criticalpoint}{T_\criticalpoint}.
    \label{eq:sm-landau-ginzburg}
\end{equation}
这里$a_0,u,\kappa>0$，系数按单位体积定义，$h$为与$\phi$共轭的外场。若$\phi=M$，上节使用的场便是$h=B$。正的梯度项意味着相邻区域的序参量取不同值要付出自由能，变化越急，代价越高。每个场构型在热平衡中的权重正比于$\exp[-\mathcal F[\phi]/(\kB T)]$。

在零场下选取一个均匀极小值$\bar\phi$，令$\delta\phi=\phi-\bar\phi$。先在平均场临界温度附近作展开，短波涨落对临界温度的移动可以在重新确定临界点时计入。把自由能展开到涨落的二次项，其局域曲率为
\[
    R=a_0t+3u\bar\phi^2=
    \begin{cases}a_0t,& t>0,\\ 2a_0|t|,& t<0.\end{cases}
\]
取体积$V$并用周期边界条件，定义$\delta\phi_{\bm q}=\int_V \d^d x\,\e^{-\i\bm q\cdot\bm x}\delta\phi(\bm x)$。此时不同波矢的涨落在高斯近似下解耦，
\[
    \mathcal F_2=\frac1{2V}\sum_{\bm q}(R+\kappa q^2)|\delta\phi_{\bm q}|^2 .
\]
对每个独立的实模式作高斯平均，就得到连通关联函数$C(\bm x)=\ave{\phi(\bm x)\phi(\bm0)}-\ave{\phi}^2$的傅里叶变换，即静态结构因子
\begin{equation}
    S(\bm q)=\int \d^d x\,\e^{-\i\bm q\cdot\bm x}C(\bm x)
    =\frac{\ave{|\delta\phi_{\bm q}|^2}}{V}
    =\frac{\kB T}{R+\kappa q^2},\qquad
    \xi=\sqrt{\frac{\kappa}{R}} .
    \label{eq:sm-gaussian-correlation}
\end{equation}
在三维，对上式作逆变换得到$C(r)=\kB T\,\e^{-r/\xi}/(4\pi\kappa r)$，其中$r=|\bm x|$大于微观尺度。关联长度$\xi$可以从实空间的衰减读出，也可以从散射实验测量的峰宽读出。散射强度降到峰值的一半时，波矢正好是$q=1/\xi$。越接近临界点，低波矢的峰越高、越窄。高斯近似给出$\xi\propto|t|^{-1/2}$。在\ssecref{sec:ising-chain}的一维格点模型中，可以直接从自旋之间的关联求出$\xi$，在那里关联长度只在零温发散。

\paragraph{Ginzburg判据}
在有序相选定一个分支，比较一个关联区域内的平均序参量$\phi_\xi$的涨落与$\bar\phi^2=a_0|t|/u$。取线度$\xi$的平滑空间平均，参与涨落的主要是$q\lesssim 1/\xi$的模式。由\eqref{eq:sm-gaussian-correlation}，忽略仅与平均窗口形状和维数有关的常数，
\[
    \ave{(\delta\phi_\xi)^2}\sim
    \int_{q\lesssim1/\xi}\frac{\d^d q}{(2\pi)^d}
         \frac{\kB T}{R+\kappa q^2}
    \sim\frac{\kB T}{\kappa}\,\xi^{2-d}.
\]
用$R=2u\bar\phi^2=\kappa/\xi^2$消去$\bar\phi$，得到平均场近似适用的Ginzburg判据
\begin{equation}
    \frac{\ave{(\delta\phi_\xi)^2}}{\bar\phi^2}
    \sim \frac{u\kB T}{\kappa^2}\,\xi^{4-d}\ll 1 .
\end{equation}
在$d>4$时，关联长度越大，这个比值越小。在$d<4$时，趋近有限温度连续相变会使比值增大，微弱的四次相互作用也会改变长波涨落。因而短程标量$\phi^4$模型的上临界维数为$4$。在$d=4$时，比值的幂次为零，需进一步考察耦合随尺度的变化，临界幂律可以带有对数修正。

\paragraph{粗粒化与普适性}
为什么液气临界点附近的密度涨落和单轴铁磁体的磁化涨落能有相同的临界指数？设$\Lambda$是空间粗粒化允许的最大波数。选取尺度因子$b>1$，先对$\Lambda/b<q<\Lambda$的短波涨落取热平均，再用$\bm x'=\bm x/b$恢复原来的波数截止，并重新标定场的大小。短波自由度的作用留在新的自由能系数中。反复进行这个变换，就是考察不同观察尺度下的有效理论。\footnote{参见David Tong，\href{https://davidtong.org/teaching/statistical-field-theory/}{\textit{Statistical Field Theory}}，第2、3章。}

在临界高斯理论中，令$\phi(\bm x)=b^{-(d-2)/2}\phi'(\bm x')$，可保持$\int \d^d x\,|\nabla\phi|^2$不变。于是二次项系数在一次变换下乘以$b^2$，四次项系数乘以$b^{4-d}$。这从另一个角度给出了上临界维数。所有无量纲耦合在变换下保持不动的理论称为固定点。固定点附近，扰动若随放大尺度而增长，称为相关扰动，若衰减则称为无关扰动。

把温度和外场调到临界点后，如果微观作用的差别沿无关方向衰减，长距离关联便由同一固定点控制。三维短程单轴铁磁体和普通单组分流体的液气临界点就属于同一个Ising普适类。流体没有精确的密度反演对称性，它的温度型和外场型标度场是温度、化学势偏离临界值的适当组合。两类系统有相同的临界指数，临界温度和关联长度的幅度仍取决于微观参数。

在短程三维Ising普适类中，临界行为由相互作用固定点控制。以$t,h$代表温度和外场方向上归一化的标度场，记它们在固定点附近的放大指数为$y_t,y_h>0$。一个线度为$b$的块含有$b^d$个原来的体积单元，自由能密度的奇异部分满足
\[
    g_{\rm sing}(t,h)=b^{-d}g_{\rm sing}(b^{y_t}t,b^{y_h}h).
\]
取$b=|t|^{-1/y_t}$，便得到上节所用的$|t|^{d/y_t}\mathcal G_\pm(h/|t|^{y_h/y_t})$。同一尺度变换给出$\xi(t,0)=b\,\xi(b^{y_t}t,0)$，从而$\nu=1/y_t$、$2-\alpha=d\nu$和$\Delta=y_h/y_t$。这里的超标度关系$2-\alpha=d\nu$适用于自由能奇异部分由单一关联长度决定的情形。在上临界维数以上，四次耦合虽然随尺度而减小，却仍决定有序相的$\bar\phi^2$，直接在标度函数中将它置零会丢失这一支。

\section{多组分系统}

我们在前文只考虑了单组分系统的情况，而许多系统都是由多个组分组成的。
当系统由$r$个组分组成时，各组分摩尔数分别为$n_1,\ldots,n_r$，
选取状态参量$T,p,n_1,\ldots,n_r$，则由Euler方程
\begin{equation}
	U=\sum_{i=1}^r\pvat{U}{n_i}{T,p,n_{\neq i}}n_i=:\sum_{i=1}^rn_iu_i.
\end{equation}
其中下标$n_{\neq i}$表示不包括组分$i$的其他组分的摩尔数，$u_i$表示在温度$T$、压强$p$和其他组分摩尔数$n_{\neq i}$不变时，组分$i$的偏摩尔内能。
同理可以定义偏摩尔体积$v_i$、偏摩尔熵$s_i$和偏摩尔Gibbs自由能$g_i$，即组分$i$的化学势：
\begin{equation}
	\mu_i=\pvat{G}{n_i}{T,p,n_{\neq i}}.
\end{equation}
内能的微分式、Euler方程和Gibbs--Duhem关系分别为
\begin{subequations}
	\begin{gather}
		\d U=T\d S-p\d V+\sum_{i=1}^r\mu_i\d n_i,\\
		U=TS-pV+\sum_{i=1}^r\mu_in_i,\\
		S\d T-V\d p+\sum_{i=1}^rn_i\d\mu_i=0.
	\end{gather}
\end{subequations}

\begin{theorem}
	{Gibbs相律}{Gibbs phase rule}
	有$r$个独立组分的系统在$t$个相共存时，系统有$T,p,\mu_1,\ldots,\mu_r$共$(r+2)$个强度量，每个相有一个Gibbs--Duhem关系。在没有其他外场或组成约束的通常共存条件下，自由度为
	\begin{equation}
		f=r+2-t.
	\end{equation}
\end{theorem}

\begin{corollary}
	对于单一组分，$r=1$，两相共存时$f=1$，在$p\vs T$相图上是一条线；三相共存时$f=0$，在$p\vs T$相图上是一个点。
\end{corollary}

\begin{example}
	{多组分理想气体}{}
	考虑$r$种可区分的理想气体，各组分温度相同。设组分$i$的定容摩尔热容为$c_iR$，它在参考温度$T_0$、摩尔体积$v_0$处的摩尔熵为$s_{i,0}$。根据\textbf{Gibbs定理}，混合物的熵等于各组分单独占据全部体积时的熵之和，
	% \footnote{虽然陈曦老师课上用统计力学证明了这个定理，但这不是基本假设之一吗？🤔}
	% 输入emoji会编译错误
	\[
		S=\sum_{i=1}^rn_i\biggb{s_{i,0}+c_iR\ln\frac{T}{T_0}+R\ln\frac{V}{n_iv_0}} .
	\]
	当气体混合时，所有气体的体积均为$V$，这相比混合前(保持$T,p$不变)的体积$V_i$增大了，由此导致了熵的增加，称为混合熵
	\begin{equation}
		\label{eqn:mixing entropy}
		S^*=-nR\sum_{i=1}^rx_i\ln x_i.
	\end{equation}
	\begin{notecaptionblock}
		\tikzfigure{mixing-entropy}
		\captionof{figure}{混合气体的过程熵会增加}
	\end{notecaptionblock}
\end{example}

\subsection{稀溶液}

考虑溶剂(solvent)中只有一种溶质(solute) B的稀溶液(solution)，
溶剂和溶质B的摩尔数分别为$N,n\,(n\ll N)$，
溶质B的摩尔分数$x:=n/(N+n)\ll 1$，
溶质B溶于溶液的混合熵
\begin{equation}
	S^*=-(N+n)R\Bigb{x\ln x+(1-x)\ln(1-x)}
	=nR(1-\ln x)+\bigo(nRx).
\end{equation}
在理想稀溶液近似下，将纯溶剂和无限稀释的溶质作为参考，可写出Gibbs自由能
\[
	G(T,p,N,n)=N\mu_0(T,p)+n\psi(T,p)-TS^*.
\]
$\mu_0$是纯溶剂的化学势，$\psi$是溶质的标准态化学势，包含溶质与溶剂相互作用的贡献。利用上式中完整的理想混合熵，在固定$T,p$时分别对$N,n$求导，得到
\begin{subequations}
	\begin{align}
		\label{eqn:mu solvent}
		\mu         & \equiv\pvat{G}{N}{T,p,n}=\mu_0+RT\ln(1-x)\simeq\mu_0-xRT,\\
		\label{eqn:mu solute}
		\mu_\mathrm{B} & \equiv\pvat{G}{n}{T,p,N}=\psi+RT\ln x.
	\end{align}
\end{subequations}

\begin{theorem}
	{渗透压与van ’t Hoff关系}{van't Hoff relation}
	以只允许溶剂通过的半透膜分隔溶液和纯溶剂。定温下阻止溶剂净流动所需的压差$\pi=p_\aqueous-p_\liquid$称为渗透压(osmotic pressure)。设$v$为纯溶剂的摩尔体积，理想稀溶液的一阶近似给出
	\begin{equation}
		\label{eqn:van't Hoff}
		\pi v\simeq xRT.
	\end{equation}
	\begin{notecaptionblock}
		\tikzfigure{osmotic-pressure}
		\captionof{figure}{半透膜引起的渗透压}
	\end{notecaptionblock}
\end{theorem}

\begin{proof}
	平衡时两端溶剂的化学势相同$\mu_0(T,p_\liquid)=\mu(T,p_\aqueous)$，代入\eqref{eqn:mu solvent}并对$p$展开：
	\[
		\mu(T,p_\aqueous)\simeq\mu_0(T,p_\aqueous)-xRT\simeq\mu_0(T,p_\liquid)+\pvat{\mu_0}pT\pi-xRT,
	\]
	由纯溶剂的$\d\mu_0=-s\d T+v\d p$即得结论。
\end{proof}

\begin{example}
	{溶液相变点随溶质的变化}{}
	没有溶质时气液共存$\mu_{\liquid,0}=\mu_\gas$；加入溶质后，溶剂的化学势变为
	\[
		\mu_\liquid=\mu_{\liquid,0}-xRT<\mu_\gas.
	\]
	如果溶质为不易挥发的物质（如盐），则气态的化学势形式不变，液相的化学势会低于气相，为了达到新的平衡需要改变温度。由$s_\liquid<s_\gas$，且
	\[
		\pv\mu{T}=-s,
	\]
	故气相化学势$\mu_\gas$的下降要快于液相，故沸点会升高，再次平衡时
	\[
		\mu_{\liquid,0}(T+\Delta T,p)-xR(T+\Delta T)=\mu_\gas(T+\Delta T,p).
	\]
	在纯溶剂的沸点处展开，保留$x$和$\Delta T$的一阶项，
	\[
		(s_\gas-s_\liquid)\Delta T\simeq xRT .
	\]
	以$\ell=T(s_\gas-s_\liquid)$表示纯溶剂的摩尔汽化潜热，沸点升高量为
	\begin{equation}
		\Delta T\simeq\frac{xRT^2}\ell .
	\end{equation}
	从液相到固相的情况正相反（因为只有液相才有溶质），熔点会降低。
\end{example}
\begin{example}
	{电化学电池}{}
	带电粒子跨相输运时，化学和静电两项共同决定其电化学势
	\[
		\widetilde\mu_i=\mu_i+z_iF\phi,\qquad F=e\NA .
	\]
	$\mu_i$是摩尔化学势，$z_i$是带符号的电荷数，$F$为Faraday常数。在等温条件下，电化学势的空间差异可驱动带电粒子输运。
	% \[J\propto-\nabla\mu=-\nabla\mu_i+ze\NA E.\]
	
	以纯金属电极的还原半反应$\mathrm{M}^{z+}+z\,\mathrm e^-\rightleftharpoons\mathrm M$为例。平衡要求参与反应的摩尔电化学势满足
	\[
		\widetilde\mu_{\mathrm{M}^{z+},\aqueous}+z\widetilde\mu_{\mathrm e,\electrode}=\mu_{\mathrm M,\electrode}.
	\]
	写$\mu_{\mathrm{M}^{z+},\aqueous}=\mu^\circ_{\mathrm{M}^{z+},\aqueous}+RT\ln a_{\mathrm{M}^{z+}}$，其中$a_{\mathrm{M}^{z+}}$是无量纲活度。固定温度、电极材料和参比电极，将组成无关的项并入标准电极电势$E^\circ$，得到还原电极电势
	\[
		E=E^\circ+\frac{RT}{zF}\ln a_{\mathrm{M}^{z+}}.
	\]
	这是电极反应的Nernst方程。\footnote{参见IUPAC \textit{Gold Book}的\href{https://goldbook.iupac.org/terms/view/09068}{Nernst equation}条目。}
\end{example}

\begin{example}
	{化学反应}{}
	考虑化学反应
	\[
		0\rightleftharpoons\sum_{i=1}^r\nu_i\mathrm{A}_i.
	\]
	以$\xi$为反应进度，$\d n_i=\nu_i\d\xi$。在恒温恒压的条件下，反应能双向进行的平衡点使Gibbs自由能取极小值，
	% \footnote{在混合系统中，总化学势一般不是各组分简单相加。}
	\[
		\d G=\sum_{i=1}^r\mu_i\d n_i=0\implies\sum_{i=1}^r\nu_i\mu_i=0.
	\]
	% 应当注意：$\mu_i$均来自于混合系统的Gibbs自由能。
	对于理想气体混合物，取标准压强$p^\circ$，组分$i$的分压为$p_i$，无量纲活度$a_i=p_i/p^\circ$。其化学势为
	\[
		\mu_i=\mu_i^\circ(T)+RT\ln a_i .
	\]
	将其代入平衡条件，定义标准反应Gibbs自由能$\Delta_{\rm r}G^\circ=\sum_i\nu_i\mu_i^\circ$，可得
	\begin{equation}
		\prod_{i=1}^r a_{i,\rm eq}^{\nu_i}=K(T)=\exp\biggl[-\frac{\Delta_{\rm r}G^\circ(T)}{RT}\biggr].
	\end{equation}
	给定标准态时，无量纲的化学平衡常数$K$只与$T$有关。若用摩尔浓度表示，可代入$p_i=(n_i/V)RT$。
\end{example}

\subsectionstar{相分离}

A和B两种组份混合后可能出现相分离，分成$\alpha$和$\beta$两相。每一相都有A和B两组份，但比例不同。%我们来看一下这是怎么发生的。假设两组份的摩尔数分别为$n_\mathrm{A}$和$n_\mathrm{B}$，总摩尔数$n=n_\mathrm{A}+n_\mathrm{B}$，
设B的摩尔比$x:=n_\mathrm{B}/n$。假设如果没有相分离，即两组份互溶，单位摩尔的Gibbs自由能为$g(x)$。
在一定温度下，$\alpha$和$\beta$两相中B的摩尔比分别为$x_\alpha$和$x_\beta$，故两相的摩尔数$n_\alpha$和$n_\beta$%下面求如果有
\begin{align*}
	\begin{cases}
		n_\alpha+n_\beta=n  \\
		x_\alpha n_\alpha+x_\beta n_\beta=xn
	\end{cases}
	\implies
	\begin{cases}
		n_\alpha=\frac{x_\beta-x}{x_\beta-x_\alpha}n  \\
		n_\beta =\frac{x-x_\alpha}{x_\beta-x_\alpha}n
	\end{cases}
\end{align*}
两相的自由能$g_\alpha=g(x_\alpha),\,g_\beta=g(x_\beta)$。若有相分离，其自由能是一条直线
\[
	g^\star(x)=\frac{n_\alpha}n g_\alpha+\frac{n_\beta}n g_\beta
	=\frac{g_\beta-g_\alpha}{x_\beta-x_\alpha}(x-x_\alpha)+g_\alpha.
\]
若$x_\alpha<x<x_\beta$时有$g^\star<g$，分成两相即可降低自由能。平衡的$x_\alpha,x_\beta$使这条直线成为$g(x)$的下公切线。

\begin{notecaptionblock}
	\tikzfigure{phase-separation}
	\captionof{figure}{相分离}
	\label{fig:phase separation}
\end{notecaptionblock}

\begin{example}
	{}{}
	下面来看怎样才能产生\figref{fig:phase separation} 所示的自由能。
	\begin{notecaptionblock}
		\tikzfigure{mixture-internal-energy}
		\tikzfigure{mixture-entropy}
		\captionof{figure}{混合物的内能和熵}
		\label{fig:mixture-energy-entropy}
	\end{notecaptionblock}
	混合物的摩尔Gibbs自由能为
	\[
		g=h-Ts=u+pv-Ts.
	\]
	在混合体积变化可忽略的模型中，混合焓可用混合内能近似。\figref{fig:mixture-energy-entropy} 给出了混合内能为正的例子，摩尔混合熵为
	\[
		s^*=-R\Bigb{x\ln x+(1-x)\ln(1-x)}.
	\]
	因此，在低温时，能量占优，可能会出现相分离
	\begin{notecaptionblock}
		\tikzfigure{mixture-gibbs-energy}
		\captionof{figure}{低温时的相分离}
	\end{notecaptionblock}
\end{example}
